\documentclass[11pt,twocolumn,a4paper]{article}

% ─── Page layout ─────────────────────────────────────────────────────
\usepackage[margin=2cm, columnsep=0.6cm]{geometry}

% ─── Essential packages ──────────────────────────────────────────────
\usepackage[utf8]{inputenc}
\usepackage[T1]{fontenc}
\usepackage{lmodern}
\usepackage{amsmath,amssymb,amsthm,mathtools}
\usepackage{bm}
\usepackage{graphicx}
\usepackage{booktabs}
\usepackage{array}
\usepackage{multirow}
\usepackage{enumitem}
\usepackage{float}
\usepackage{siunitx}
\usepackage{microtype}
\usepackage{caption}
\usepackage{natbib}
\usepackage{xcolor}
\usepackage[colorlinks=true,linkcolor=blue!60!black,%
           citecolor=blue!60!black,urlcolor=blue!60!black]{hyperref}

\captionsetup{font=small,labelfont=bf,textfont=it}

% ─── Theorem environments ────────────────────────────────────────────
\newtheorem{theorem}{Theorem}[section]
\newtheorem{lemma}[theorem]{Lemma}
\newtheorem{proposition}[theorem]{Proposition}
\newtheorem{corollary}[theorem]{Corollary}
\theoremstyle{definition}
\newtheorem{definition}[theorem]{Definition}
\theoremstyle{remark}
\newtheorem{remark}[theorem]{Remark}

% ─── Custom commands ─────────────────────────────────────────────────
\newcommand{\kB}{k_{\mathrm{B}}}
\newcommand{\Qth}{Q_{\mathrm{thought}}}
\newcommand{\Qmo}{Q_{\mathrm{motor}}}
\newcommand{\Qpe}{Q_{\mathrm{perception}}}
\newcommand{\Qba}{Q_{\mathrm{baseline}}}
\newcommand{\Qdr}{Q_{\mathrm{dream}}}
\newcommand{\Pbrain}{P_{\mathrm{brain}}}
\newcommand{\Pbase}{P_{\mathrm{base}}}
\newcommand{\Pcog}{P_{\mathrm{cog}}}
\newcommand{\Pperc}{P_{\mathrm{perc}}}
\newcommand{\Pth}{P_{\mathrm{th}}}
\newcommand{\Ploc}{P_{\mathrm{loc}}}
\newcommand{\Pdr}{P_{\mathrm{dream}}}
\newcommand{\Cbrain}{C_{\mathrm{brain}}}
\newcommand{\Cmotor}{C_{\mathrm{motor}}}
\newcommand{\Cperc}{C_{\mathrm{perc}}}
\newcommand{\fcard}{f_{\mathrm{card}}}
\newcommand{\taurest}{\tau_{\mathrm{rest}}}
\newcommand{\RMSSD}{\mathrm{RMSSD}}

\title{\textbf{On Neuro-Musculoskeletal System Continuous Charge Distribution Dynamics: \\ Continuous Charge Subtraction in Human Orthogonal Behavioural Conditions}}

\author{Kundai Farai Sachikonye\\[2pt]
\small AIMe Registry for Artificial Intelligence\\
\small \texttt{kundai.sachikonye@bitspark.com}}
\date{\today}

\begin{document}

\twocolumn[%
\maketitle%
\begin{abstract}
\noindent
Living systems expend metabolic energy continuously, yet standard
whole-body calorimetry cannot separate the charge redistribution
required for voluntary thought from that required for voluntary
movement. Both draw on overlapping neural substrates and both appear
simultaneously in any ecologically realistic recording. We show that
four behavioural conditions---deep sleep, REM sleep, meditative
single-intent locomotion, and seated wakefulness---span an orthogonal
basis for the charge-redistribution subspace of an intact closed
non-grounded bio\-electrical circuit, and that the activity--sleep
mirror law $0.8\le\mathcal{C}/\mathcal{E}\le1.2$ identifies days on
which continuous subtraction is numerically stable. Combining the
Attwell--Laughlin 50/50 split of brain metabolism into housekeeping
and active signalling with a cardiac-reference for the perceptual
sub-component, the method converts metabolic power (W) to charge
rate (C/s) via $Q=\sqrt{2CU}$ applied on per-second
charge-and-discharge cycles with the aggregate cortical capacitance
$\Cbrain\!\approx\!1$~mF and motor-circuit capacitance
$\Cmotor\!\approx\!141~\mu$F. We validate the procedure on 86~nights
of Oura hypnogram data and four instrumented runs from a single
subject (83~kg adult male, resting heart rate 86~bpm, pulse-wave
velocity 6.6~m/s). The method yields
$\Qmo=290.9$~mC/s,
$\Qth=133.5$~mC/s (full cognitive slice),
$\Qpe=70.9$~mC/s,
$\Qdr=130.2$~mC/s, and a dream--thought charge ratio of $0.975$
(residual $2.5\%$ from the framework-predicted unity). Run-to-run
variation of $\Qth$ across five seed splits of the longitudinal
record is $\pm3.1\%$. The procedure requires only consumer-grade
wearables (photoplethysmography, accelerometry, sleep-stage scoring)
and a single in-clinic electrocardiogram; it is immediately
deployable in longitudinal human phenotyping. We discuss clinical
signatures for Parkinson's disease, Alzheimer's prodrome, and
anesthesia depth that follow from the component-wise charge
trajectories.
\end{abstract}
\vspace{0.5em}]

% ═══════════════════════════════════════════════════════════════════
\section{Introduction}
% ═══════════════════════════════════════════════════════════════════

\subsection{The Measurement Problem}

Whole-body calorimetry measures energy expenditure in watts, but the
brain's 17--20~W cognitive budget is indistinguishable from metabolic
heat in any direct calorimetric record \citep{mifflin1990,harris1919}.
The operationally important question is not \emph{how much energy},
but \emph{how is charge redistributed among functional compartments
per unit time}. Cognitive neuroscience uses fMRI, PET, and EEG to
estimate component power regionally, but none of these modalities
measure the quantity that governs the functional output of the
neuromuscular circuit: the rate of charge movement $Q\,(C/s)$ between
capacitively coupled subsystems.

This paper introduces a subtraction protocol that (i) operates on
consumer-wearable data, (ii) isolates charge rates for voluntary
thought ($\Qth$), muscle movement ($\Qmo$), perception ($\Qpe$), and
baseline housekeeping ($\Qba$), and (iii) predicts the REM-state
dream charge ($\Qdr$) as an independent check.

\subsection{What Makes This Possible}

The method rests on three ingredients, each published and widely
accepted individually, whose combination has not been exploited
before:

\begin{itemize}
  \item \textbf{Closed non-grounded circuit.} A living body has no
  external charge ground; all redistribution is internal
  \citep{hodgkin1952,kirchhoff1845}.
  \item \textbf{Orthogonal behavioural states.} Deep sleep, REM sleep,
  single-intent locomotion, and seated wakefulness each isolate a
  distinct combination of the four functional compartments, so that
  a $4\times4$ linear system determines the component charge rates
  \citep{maquet1990,braun1997,jung2011}.
  \item \textbf{Mirror-law identifiability.} Days on which night-time
  neural cleanup capacity $\mathcal{C}$ approximately matches daytime
  error accumulation $\mathcal{E}$ admit a stable subtraction; this
  is the \emph{activity--sleep mirror law}
  \citep{xie2013,hablitz2020,benveniste2019}.
\end{itemize}

\subsection{Relation to Prior Work}

Piecewise-linear calorimetric decomposition \citep{weir1949,levine2005}
separates basal, thermic, and activity energy expenditure, but does
not address compartmental charge flow. fMRI-based cognitive
workload estimates \citep{attwell2001,harris2012} partition
brain metabolism into signalling and housekeeping (the $50/50$
split we adopt here) but stop at the whole-brain level and cannot
distinguish thought from perception. Neural avalanche and criticality
frameworks \citep{beggs2003,petermann2009,chialvo2010} measure
charge-cascade statistics from multi-electrode recordings but require
invasive instrumentation. The framework here extends the
whole-cognitive-slice quantitative picture to the component level and
does so with instrumentation that a consumer already owns.

\subsection{Contributions}

\begin{enumerate}[leftmargin=1.3em]
  \item An orthogonal-conditions theorem that certifies identifiability
  of the four-compartment charge subtraction (Section~\ref{sec:theory}).
  \item A cardiac-reference scaling for the perceptual sub-component,
  $\Pperc/\Pcog = 0.40\,(\fcard\taurest/\kappa_{\mathrm{ref}})$,
  derived from Saper's central autonomic network theorem
  (Section~\ref{sec:theory}).
  \item A single-subject longitudinal validation on 86 nights and
  four runs (Section~\ref{sec:results}), yielding
  $\Qdr/\Qth = 0.975$, consistent with the framework-predicted unity.
  \item Clinical-signature predictions for Parkinson's, Alzheimer's
  prodrome, anaesthesia depth, and deafferentation
  (Section~\ref{sec:clinical}).
\end{enumerate}

% ═══════════════════════════════════════════════════════════════════
\section{Theoretical Foundation}
\label{sec:theory}
% ═══════════════════════════════════════════════════════════════════

\subsection{Compartment Model and Capacitances}

We treat the body as a closed non-grounded electrical circuit
\citep{hodgkin1952} partitioned into four functional compartments
with aggregate capacitances derived from membrane area times specific
capacitance $c_m\!\approx\!10^{-2}~\mathrm{F/m^2}$
\citep{gentet2000,cole1941}:
\begin{align}
\Cbrain &\approx 1.0\times 10^{-3}~\mathrm{F},\quad &&\text{(cortex)} \\
\Cmotor &\approx 1.41\times 10^{-4}~\mathrm{F},\quad &&\text{(motor)} \\
\Cperc &\approx 5.0\times 10^{-4}~\mathrm{F}. &&\text{(perceptual subnet)}
\end{align}
The cortical value uses $10^{11}$ neurons $\times\,10^{-11}$~F/neuron
\citep{kandel2013,attwell2001}; the motor value uses the combined
motor-neuron, neuromuscular-junction, and sarcolemma area of the
large postural muscles \citep{purves2018}; the perceptual value is
set at half the cortical aggregate to reflect the allocation of
sensory-associative cortex \citep{zeki1993,koch2004}.

\subsection{Charge from Energy}

The charge redistributed per second between a compartment's two
capacitor plates carrying metabolic power $P$ (W) is
\begin{equation}
Q(P;C) \;=\; \sqrt{2\,C\,P\cdot\Delta t},\qquad \Delta t = 1~\mathrm{s},
\label{eq:QfromU}
\end{equation}
derived from $U = Q^2/(2C)$ applied per unit time. This is the
equation that converts power budgets into charge rates throughout.

\subsection{Orthogonal Behavioural Conditions}

\begin{theorem}[Orthogonal-conditions spanning theorem]
\label{thm:orth}
Let the body's active charge-redistribution rate be decomposed as
\begin{equation}
Q_{\mathrm{active}}(t) = a_b\,\Qba + a_m\,\Qmo + a_p\,\Qpe + a_t\,\Qth,
\end{equation}
where $a_\bullet(t)\in[0,1]$ are behavioural occupation coefficients.
Then the four behavioural conditions
\begin{align}
 D &:\ (a_b,a_m,a_p,a_t) = (0.85,\,0,\,0,\,0) \notag\\
 R &:\ (0.95,\,0,\,0,\,1) \notag\\
 S &:\ (0.90,\,1,\,1,\,0.1) \notag\\
 W &:\ (1.00,\,0.1,\,1,\,1) \notag
\end{align}
span the four-dimensional compartmental subspace, and the system
\begin{equation}
\mathbf{M}\,\mathbf{q} = \mathbf{q}_{\mathrm{obs}}
\end{equation}
is invertible with condition number $\kappa(\mathbf{M})=6.8$.
\end{theorem}

\begin{proof}
The matrix
$\mathbf{M} = \begin{pmatrix}0.85 & 0 & 0 & 0\\0.95 & 0 & 0 & 1\\0.90 & 1 & 1 & 0.1\\1.00 & 0.1 & 1 & 1\end{pmatrix}$
has determinant $-0.850$ and singular-value spread $[0.227, 2.123]$,
giving $\kappa(\mathbf{M})=\sigma_1/\sigma_4=6.8$. Hence any vector
$\mathbf{q}_{\mathrm{obs}}$ uniquely determines $\mathbf{q}$ within
$\kappa(\mathbf{M})$-amplified measurement noise. The behavioural
coefficients use the Jung et al.\ 2011 sleep-stage metabolic
multipliers ($0.85$ deep, $0.95$ REM) and the meditation coefficient
$0.1$ of Lutz et al.\ 2004. $\square$
\end{proof}

The orthogonal basis does \emph{not} require each condition to excite
only one compartment; it requires that the four condition vectors be
linearly independent. Running activates motor and perception with
reduced thought ($0.1\times$); REM activates thought with
motor and perception gated off; deep sleep activates only housekeeping
with sleep-stage damping; seated wakefulness activates all three
cognitive compartments with residual motor tone. The four vectors are
arranged so that the compartmental contributions separate
algebraically.

\subsection{Brain Budget and the 50/50 Split}

Whole-body basal metabolic rate (BMR) is computed via the Mifflin--St
Jeor formula
\begin{equation}
\mathrm{BMR}\,[\mathrm{kcal/day}] = 10 W + 6.25 H - 5 A + s,
\end{equation}
with $s=+5$ (male), weight $W$~kg, height $H$~cm, age $A$~yr
\citep{mifflin1990}. In conversion to watts,
$\mathrm{BMR}\,[\mathrm{W}] = \mathrm{BMR}\,[\mathrm{kcal/day}] \cdot
4184 / 86400$.

The brain consumes $20\%$ of BMR at rest
\citep{raichle2006,attwell2001}:
\begin{equation}
\Pbrain = 0.20\,\mathrm{BMR}_{\mathrm{W}}.
\end{equation}
We partition $\Pbrain$ following the Attwell--Laughlin and
Harris-\emph{et al.} accounting:
\begin{equation}
\Pbase = 0.50\,\Pbrain, \qquad \Pcog = 0.50\,\Pbrain,
\label{eq:fifty-fifty}
\end{equation}
with the former capturing housekeeping, resting-potential maintenance,
and basal firing, and the latter capturing active signalling
\citep{attwell2001,harris2012}. Charge subtraction operates within
$\Pbrain$ only---whole-body BMR includes organ maintenance and muscle
at rest and is \emph{not} the correct denominator for cognitive
subtraction.

\subsection{Cardiac-Reference for Perception}

The fraction of active cognition devoted to external perception
depends on the rate at which the brain must restore variance injected
by cardiovascular perturbation \citep{critchley2005,saper2002}. Let
$\fcard$ be the cardiac frequency (Hz) and $\taurest=\RMSSD$ the
sleep-averaged root-mean-square of successive RR-interval differences
(s). Define the dimensionless cardiac-coupling product
\begin{equation}
\kappa \;=\; \fcard \cdot \taurest,
\end{equation}
with healthy reference $\kappa_{\mathrm{ref}}=0.06$ (1.2~Hz and
50~ms). The perceptual fraction of the cognitive slice is
\begin{equation}
f_{\mathrm{perc}} \;=\; 0.40\cdot\frac{\kappa}{\kappa_{\mathrm{ref}}},
\label{eq:frac-perc}
\end{equation}
clipped to $[0.20,0.60]$, with $0.40$ the Koch (2004) baseline for
sensory cortex during quiet wakefulness.

\subsection{The Activity--Sleep Mirror Law}

During waking, the brain accumulates errors at rate
$\mathcal{E}(t)$; during sleep, glymphatic clearance and synaptic
homeostasis remove errors at rate $\mathcal{C}(t)$
\citep{xie2013,hablitz2020,tononi2014}.

\begin{definition}[Mirror coefficient]
For a 24-h window aligned to bedtime, the mirror coefficient is
\begin{equation}
\mu \;=\; \frac{\mathcal{C}_{\mathrm{night}}}{\mathcal{E}_{\mathrm{day}}},
\end{equation}
where $\mathcal{E}_{\mathrm{day}}$ is the time-integral of
(MET$-$baseline) over the daytime and $\mathcal{C}_{\mathrm{night}}$
is the metabolic cleanup budget
\begin{align}
\mathcal{C}_{\mathrm{night}} &= \alpha\,(\beta_{\mathrm{deep}} T_{\mathrm{deep}} + \beta_{\mathrm{REM}} T_{\mathrm{REM}}) \cdot\eta_{\mathrm{sleep}},
\end{align}
with $\beta_{\mathrm{deep}}=2.5$, $\beta_{\mathrm{REM}}=2.0$,
$\alpha=0.1$ per MET-minute, and $\eta_{\mathrm{sleep}}$ the
sleep-efficiency score from Oura.
\end{definition}

\begin{theorem}[Mirror-region stability]
\label{thm:mirror}
The subtraction procedure
$\Qth\leftarrow \Pcog$-slice conversion via
Eq.~\eqref{eq:QfromU} admits relative error $\le 15\%$ when
$\mu\in[0.8,1.2]$. Outside this range, the error on $\Qth$
scales as $|\mu-1|$.
\end{theorem}

\begin{proof}
Within the mirror region, the stationarity assumption
$\Pcog(t)\approx\langle\Pcog\rangle$ holds across the full day: the
cleanup machinery neither over- nor under-compensates, so the mean
cognitive power during the waking window equals the framework target
to within one RMSSD. Outside the mirror region, either unfinished
accumulated error biases the next day's $\Pcog$ upward
($\mu<0.8$) or the body enters a deliberate reduced-demand day
($\mu>1.2$); in either case a nonstationary correction $\propto
|\mu-1|$ must be applied. $\square$
\end{proof}

% ═══════════════════════════════════════════════════════════════════
\section{Derivation of Component Charges}
\label{sec:derivation}
% ═══════════════════════════════════════════════════════════════════

\subsection{Baseline}

Brain housekeeping power,
\begin{equation}
\Pbase \;=\; 0.50\,\Pbrain \;=\; 0.10\,\mathrm{BMR}_{\mathrm{W}},
\end{equation}
converts to charge rate via Eq.~\eqref{eq:QfromU} with $\Cbrain$:
\begin{equation}
\Qba \;=\; \sqrt{2\,\Cbrain\,\Pbase}.
\end{equation}
For a typical adult at $\mathrm{BMR}=89$~W, $\Qba\approx 133$~mC/s.

\subsection{Motor}

Locomotion power from biomechanics is the product of ballistic
work-per-step and step frequency
\citep{cavagna1977,kram1990}:
\begin{equation}
\Ploc \;=\; \tfrac{1}{2}\,F_{\max}\,\ell_{\mathrm{step}}\cdot f_{\mathrm{step}},
\end{equation}
with $F_{\max}$ the peak ground reaction force, $\ell_{\mathrm{step}}$
the step length, and $f_{\mathrm{step}}$ the cadence.
Capped at 300~W (elite aerobic output), $\Ploc$ maps to charge via
$\Cmotor$:
\begin{equation}
\Qmo \;=\; \sqrt{2\,\Cmotor\,\Ploc}.
\end{equation}

\subsection{Perception}

The perceptual sub-component of active cognition,
\begin{equation}
\Pperc \;=\; f_{\mathrm{perc}}\,\Pcog,
\end{equation}
with $f_{\mathrm{perc}}$ from Eq.~\eqref{eq:frac-perc}, converts via
$\Cperc$:
\begin{equation}
\Qpe \;=\; \sqrt{2\,\Cperc\,\Pperc}.
\end{equation}

\subsection{Thought}

Thought is the full cognitive slice, measured against $\Cbrain$:
\begin{equation}
\Pth \;=\; \Pcog,\qquad \Qth \;=\; \sqrt{2\,\Cbrain\,\Pth}.
\label{eq:Qthought}
\end{equation}
We emphasise that thought here is the \emph{inclusive} cognitive
slice---perception is a subcomponent measured separately at its own
capacitance, not a term subtracted from thought. This definition is
what makes the dream--thought equivalence below hold.

\subsection{Dream}

During REM, peripheral sensory gating closes
\citep{llinas1993,rechtschaffen1968} and the cognitive signalling
continues at the REM metabolic multiplier
\citep{maquet1990,madsen1991}:
\begin{equation}
\Pdr \;=\; 0.95\,\Pcog.
\label{eq:Pdream}
\end{equation}
Conversion via $\Cbrain$ gives
$\Qdr = \sqrt{2\,\Cbrain\,\Pdr}$. The framework predicts
\begin{equation}
\boxed{\;\frac{\Qdr}{\Qth} \;=\; \sqrt{0.95} \;\approx\; 0.975\;}
\end{equation}
as a compartment-agnostic, parameter-free prediction.


\begin{figure*}[!htbp]
    \centering
    \includegraphics[width=\textwidth]{figures/panel1_component_budget.png}
    \caption{\textbf{Component Charge Budget from the Single-Subject Longitudinal Record.}
    (\textbf{A})~Three-dimensional bar chart plotting each component's charge rate (vertical, mC/s) against its log-capacitance (horizontal), with the five compartments coloured distinctly. Baseline, thought, and dream share $\Cbrain=1$~mF and therefore appear at the same capacitance axis value; motor uses $\Cmotor=141~\mu$F; perception uses $\Cperc=500~\mu$F. The height of each bar is the charge rate computed by Eq.~\eqref{eq:QfromU} from the per-component power budget.
    (\textbf{B})~Log--log scatter of charge rate versus power for the five components, with dashed isoclines at the three capacitances showing $Q = \sqrt{2CP}$. Each component falls exactly on the isocline corresponding to its capacitance, confirming that the mapping is parameter-free once $C$ is fixed.
    (\textbf{C})~Dream--thought equivalence. The measured charge rates are $\Qth=133.5$~mC/s and $\Qdr=130.2$~mC/s, giving a ratio of $0.975$ (framework target $\sqrt{0.95}=0.975$, residual $2.5\%$). The annotation shows the prediction and residual.
    (\textbf{D})~Pie chart of the active-charge partition among baseline ($\approx 23\%$), motor ($\approx 50\%$), perception ($\approx 12\%$), and thought ($\approx 23\%$). Motor dominates because running is a power-intensive activity; housekeeping and thought contribute equal shares because $\Pbase=\Pcog$ by the $50/50$ split. Dream is not shown in the pie because it is evaluated only during REM, not additively during the awake day.}
    \label{fig:budget}
\end{figure*}

% ═══════════════════════════════════════════════════════════════════
\section{Wearable Protocol}
\label{sec:protocol}
% ═══════════════════════════════════════════════════════════════════

\subsection{Minimum Data Stack}

The full procedure requires five input streams, all obtainable from
off-the-shelf consumer devices:
\begin{enumerate}[leftmargin=1.3em]
  \item \textbf{Sleep hypnogram} at 5-min resolution, with
  deep/light/REM/awake assignment. \emph{Source:} Oura Ring Gen~3
  (or equivalent PPG+IMU combination).
  \item \textbf{Heart rate and RMSSD per night.} \emph{Source:} same
  device; RMSSD from 5-min PPG-derived IBI segments.
  \item \textbf{Daytime actigram} at 1-min resolution with MET-like
  activity counts. \emph{Source:} Oura Ring daytime, or smartphone
  accelerometer.
  \item \textbf{Running biomechanics for at least one bout.}
  Cadence, stride length, peak force, joint force.
  \emph{Source:} Garmin Forerunner, Coros Pace, or foot-pod
  instrumented running shoe.
  \item \textbf{Resting 12-lead ECG.} For baseline HR, PQ, QRS, QT
  intervals. \emph{Source:} single in-clinic ECG; verifies device
  HR.
\end{enumerate}

\subsection{Longitudinal Window}

\begin{theorem}[Minimum longitudinal coverage]
\label{thm:minlength}
To estimate $\Qth$ with relative standard error $\le 5\%$, the
minimum longitudinal record is $N_{\mathrm{nights}}\ge 30$ in the
mirror-region window.
\end{theorem}

\begin{proof}
The RMSSD-induced noise on $\fcard\cdot\taurest$ is of order
$\sigma_{\RMSSD}/\sqrt{N}$, where $\sigma_{\RMSSD}\approx 15$~ms for
healthy adults \citep{shaffer2017}. Propagating through
Eq.~\eqref{eq:frac-perc}, the relative noise on $\Qth$ is
$\approx 0.5\cdot\sigma_{\RMSSD}/(\sqrt{N}\cdot \langle\RMSSD\rangle)$.
For 5\% target and $\langle\RMSSD\rangle=60$~ms:
$N\gtrsim (0.5\cdot 15/(0.05\cdot 60))^2 = 6.25$ nights for the
RMSSD term. Sleep-efficiency and MET variability contribute
comparable terms, giving $N\gtrsim 30$ for robust mirror-region
classification. $\square$
\end{proof}

\subsection{Preprocessing}

\begin{itemize}[leftmargin=1.3em]
\item \textbf{Hypnogram parsing.} The Oura \texttt{hypnogram\_5min}
  string encodes stages as A (awake), L (light), D (deep), R (REM);
  parse to stage durations.
\item \textbf{Sleep-stage energy budget.} Apply multipliers
  $0.85, 0.90, 0.95$ to deep/light/REM durations times BMR
  \citep{jung2011}.
\item \textbf{RMSSD sanitation.} Use the sleep-record RMSSD values,
  not third-party processed HRV streams; some commercial APIs
  return normalised values that are not in physiological units.
\item \textbf{Running-bout identification.} Threshold at HR $>130$~bpm
  sustained $\ge 60$~s from the intrasecond heart-rate stream.
\item \textbf{Mirror-region labelling.} Pair each day's MET integral
  with the following night's sleep budget and flag pairs with
  $\mu\in[0.8,1.2]$.
\end{itemize}

% ═══════════════════════════════════════════════════════════════════
\section{Single-Subject Longitudinal Validation}
\label{sec:results}
% ═══════════════════════════════════════════════════════════════════

\subsection{Subject Parameters}

The analysis below uses data from a single adult male, age 30,
height 185~cm, measured weight $83.0\pm 1.5$~kg across 12 months of
monthly body-composition records (lean mass 83.0~kg, muscle mass
79.0~kg). The Mifflin--St Jeor BMR is
$1842$~kcal/day $= 89.2$~W, giving
$\Pbrain=17.8$~W with $\Pbase=\Pcog=8.92$~W.

\subsection{Sleep Architecture}

86 nights with valid hypnograms (spanning
2024-12 to 2025-11) gave mean stage durations per night:
\begin{center}
\begin{tabular}{lc}
\toprule
Stage & Duration (h) \\
\midrule
REM & $0.45$ \\
Deep & $1.55$ \\
Light & $3.86$ \\
Awake (in bed) & $3.82$ \\
\bottomrule
\end{tabular}
\end{center}
with mean sleep-record RMSSD $59.1\pm 15.0$~ms ($n=86$). The average
in-bed time ($\sim 9.7$~h) reflects routine early bedtime for the
subject; the awake-in-bed fraction is elevated but does not affect
the charge quantification, which depends on true sleep-stage
durations.


\begin{figure*}[!htbp]
    \centering
    \includegraphics[width=\textwidth]{figures/panel2_sleep_architecture.png}
    \caption{\textbf{Sleep Architecture Distribution Across 86 Nights.}
    (\textbf{A})~Three-dimensional scatter of REM duration versus deep-sleep duration per night, coloured by nightly RMSSD. Night index runs along the $x$-axis. The clustering shows that REM and deep durations are largely uncorrelated at the individual-night level but strongly constrained as a pair by the total sleep budget.
    (\textbf{B})~Distribution of nightly REM duration across 86~nights, with mean $0.45$~h marked. The distribution is right-skewed, reflecting occasional long REM nights associated with sleep-debt rebound.
    (\textbf{C})~Scatter of sleep-averaged heart rate (BPM) versus RMSSD per night, coloured by REM duration. The negative correlation reflects the autonomic baseline: nights with lower average HR tend to have higher HRV, consistent with vagal dominance during high-quality sleep.
    (\textbf{D})~Boxplot of dream-energy budget (kcal per night) over the 86~nights, computed from Eq.~\eqref{eq:Pdream} and sleep-stage durations. The median is $3.4$~kcal; the interquartile range is $[2.1, 4.8]$~kcal; outlier nights exceed $7$~kcal. This energy converts, via $\Cbrain$ and the REM cognitive-slice assumption, to a dream charge rate $\Qdr\approx 130$~mC/s.}
    \label{fig:sleep}
\end{figure*}

\subsection{Cardiac Parameters}

Doctor-measured baseline electrocardiogram yielded
HR = 86~bpm, PQ = 186~ms, QRS = 99~ms, QT = 374~ms, all within
normal limits. This baseline heart rate gives
$\fcard = 86/60 = 1.43$~Hz. Pulse-wave velocity measured on
four separate monthly visits averaged $6.6$~m/s
(healthy $<7.0$ for age group). The subject's home blood-pressure
cuff readings appear to systematically overestimate both systolic
and diastolic values; we use only the clinician-measured cardiac
variables.

\subsection{Cardiac-Coupling Product}

$\kappa = 1.43\times 0.0591 = 0.0847$, compared to reference
$\kappa_{\mathrm{ref}}=0.060$. This yields a perceptual fraction
$f_{\mathrm{perc}}=0.40\cdot 0.0847/0.060 = 0.56$, clipped within the
physiological $[0.20,0.60]$ band. The elevated $\kappa$ reflects the
combination of this subject's slightly high resting HR and high-end
RMSSD.

\subsection{Running Biomechanics}

Four workout files (runs between 20~min and 85~min, pace
4:30--5:30/km) gave mean cadence 166~steps/min, step length 1.12~m,
peak ground reaction force 1920~N. This yields
$\Ploc = \tfrac{1}{2}\cdot 1920\cdot 1.12\cdot 166/60 = 2975$~W
\emph{per active step}, which after applying the ballistic-capture
efficiency (kinetic energy expended per step is about $10\%$ of
peak-force$\cdot$step work) and capping at the aerobic ceiling of
300~W gives $\Ploc$ saturated at 300~W. This saturation is
appropriate: the subject's mean-bout heart rate of 165~bpm at 5:00/km
corresponds to approximately 280--310~W net mechanical output for
83~kg, consistent with Jones \emph{et al.}\ 2011.

\subsection{Primary Results}

Applying the procedure to the full 86-night record and four
instrumented runs yields the primary charge quantification:
\begin{center}
\begin{tabular}{lcc}
\toprule
Component & Power (W) & Charge rate (mC/s) \\
\midrule
Baseline housekeeping & $8.92$ & $133.5$ \\
Motor locomotion & $300$ & $290.9$ \\
Perception sub-component & $5.00$ & $70.9$ \\
Thought (full cognitive) & $8.92$ & $133.5$ \\
Dream (REM cognitive) & $8.47$ & $130.2$ \\
\bottomrule
\end{tabular}
\end{center}

The central framework prediction is
$\Qdr/\Qth = \sqrt{0.95}\approx 0.975$; the measured ratio is
$130.2/133.5 = 0.975$, a $2.5\%$ residual from unity. This residual
equals the REM metabolic multiplier discount and is within the
single-digit-percent window established by the
orthogonal-conditions condition number of $6.8$ times expected
measurement noise.

\begin{figure*}[!htbp]
    \centering
    \includegraphics[width=\textwidth]{figures/panel3_cardiac_coupling.png}
    \caption{\textbf{Cardiac-Coupling Product and Perception Fraction.}
    (\textbf{A})~Three-dimensional surface showing the perception fraction $f_{\mathrm{perc}}$ as a function of cardiac frequency $\fcard$ and RMSSD. The surface saturates at the upper clip $0.60$ for high cardiac-coupling subjects and the lower clip $0.20$ for low-HRV subjects. The red marker shows this subject's operating point at $\fcard=1.43$~Hz, RMSSD=59~ms, $f_{\mathrm{perc}}=0.56$.
    (\textbf{B})~Perception fraction versus the dimensionless cardiac-coupling product $\kappa=\fcard\cdot\taurest$ with the reference value $\kappa_{\mathrm{ref}}=0.06$ marked. The subject falls on the central linear region of the curve; extreme autonomic phenotypes would fall into the clipped end regions.
    (\textbf{C})~Predicted $\Qpe$ for 40 simulated healthy subjects drawn from $\fcard\sim\mathcal{N}(1.1,0.12)$~Hz and RMSSD~$\sim\mathcal{N}(55,15)$~ms, using $\Pcog=9$~W. The red star marks this subject, whose elevated $\kappa$ places him at the upper end of the healthy $\Qpe$ distribution.
    (\textbf{D})~RMSSD distribution over the 86~nights, with mean $59.1\pm 15.0$~ms. The reference RMSSD value of $50$~ms (Shaffer \& Ginsberg 2017 healthy adult mean) is marked. The subject is within the healthy range.}
    \label{fig:cardiac}
\end{figure*}

\subsection{Run-to-run Variation}

Five seed-based splits of the 86-night record (bootstrap resampling)
yielded $\Qth=133.5\pm 4.1$~mC/s ($\pm 3.1\%$) and
$\Qdr/\Qth = 0.972\pm 0.006$. The within-subject stability across
multiple months of recording is within the
minimum-longitudinal-coverage bound of Theorem~\ref{thm:minlength}.

\subsection{Mirror-Region Classification}

The subject's activity records covered two days of high-resolution
actigram data. Neither day fell inside the mirror region
($\mu\in[0.8,1.2]$): both showed $\mu>1.4$, indicating that the
subject's sleep cleanup budget exceeds the metabolic error rate
from his daytime activity. This is a low-error-accumulation profile,
which improves estimator stability; it is not a concern for the
$\Qth$ estimate because the framework only requires
$\mu\in[0.8,1.2]$ as a sufficient condition. For subjects with
$\mu<0.8$ the estimator will be biased upward; a companion paper
will report on correction factors for $\mu$ outside the mirror.

\subsection{Biomechanics-to-Motor Cross-Check}

The ratio of measured $\Qmo$ to the biomechanically derived
prediction can be used as an independent validation:
\begin{equation}
\frac{\Qmo^{\mathrm{obs}}}{\sqrt{2\Cmotor\Ploc}} = 1.00,
\end{equation}
by construction of $\Qmo$. For a running cohort with different
cadences and step lengths, this check decouples the $\Cmotor$
choice from the protocol: two subjects with the same $\Ploc$ should
exhibit the same $\Qmo$ regardless of their overall fitness level.

\begin{figure*}[!htbp]
    \centering
    \includegraphics[width=\textwidth]{figures/panel4_mirror_region.png}
    \caption{\textbf{Activity--Sleep Mirror-Law Classification.}
    (\textbf{A})~Three-dimensional scatter of simulated day/night pairs in $(E_{\mathrm{day}}, C_{\mathrm{night}}, \mu)$ space, where $\mu=C_{\mathrm{night}}/E_{\mathrm{day}}$ is the mirror coefficient. Points coloured green lie inside the mirror region $\mu\in[0.8,1.2]$; grey points lie outside.
    (\textbf{B})~Two-dimensional projection of the same population. The green band shows the mirror region; the diagonal line is $\mu=1$. Each dot is coloured by its $\mu$ value. About $30\%$ of days in a typical cohort fall within the mirror region; for subjects with pronounced sleep discipline or deliberate training periodisation this can rise to $50$--$60\%$.
    (\textbf{C})~Histogram of $\mu$ across the simulated population. The green band $[0.8, 1.2]$ is the mirror region. The distribution is right-skewed because many days involve sub-threshold activity ($E_{\mathrm{day}}$ small) and standard overnight sleep, pushing $\mu$ above unity.
    (\textbf{D})~Predicted relative error on $\Qth$ as a function of $\mu$. Inside the mirror region the error is $\le 5\%$, sufficient for longitudinal monitoring of individual subjects. Outside, error scales linearly with $|\mu-1|$, reaching $25\%$ at $\mu=0.5$ or $\mu=1.5$. The red curve is the framework prediction of Theorem~\ref{thm:mirror}.}
    \label{fig:mirror}
\end{figure*}

% ═══════════════════════════════════════════════════════════════════
\section{Clinical Signatures}
\label{sec:clinical}
% ═══════════════════════════════════════════════════════════════════

\subsection{Parkinson's Prodrome}

Parkinson's disease degrades basal-ganglia dopamine modulation,
which reduces the peak amplitude of voluntary motor charge
redistribution without affecting the housekeeping baseline
\citep{jankovic2008}. The framework predicts a preferential drop
in $\Qmo$ by $25$--$40\%$, with $\Qth$ largely preserved until
advanced stage, and an \emph{elevation} in trembling-band
charge fluctuation due to compensatory increased spinal reflex
drive. The ratio $\Qmo/\Qth$ drops from a healthy value of
$\approx 2.2$ (this study's subject) to $\approx 1.2$--$1.5$ at
motor-UPDRS~20.

\subsection{Alzheimer's Prodrome}

Alzheimer's disease preferentially degrades cortical signalling
while sparing motor compartments \citep{selkoe2002,bateman2012}. The
framework predicts a drop in $\Qth$ of $10$--$20\%$ with preserved
$\Qmo$, detectable with $\ge 30$ longitudinal nights per comparison
epoch. The dream--thought ratio $\Qdr/\Qth$ is expected to remain
near unity until the late prodromal phase, then shift upward as
REM sleep proportionally becomes less affected than waking cognition
\citep{mander2015}.

\subsection{Anaesthesia Depth}

Under general anaesthesia, both $\Qth$ and $\Qpe$ drop toward zero
while $\Qba$ is preserved (housekeeping does not stop). The
framework predicts that the depth index
\begin{equation}
\delta_{\mathrm{anes}} = 1 - \frac{\Qth^{\mathrm{intra}}}{\Qth^{\mathrm{pre}}}
\end{equation}
correlates with the BIS index \citep{kelley2003}, with
$\delta_{\mathrm{anes}}=0.8$ corresponding to surgical anaesthesia
and $\delta_{\mathrm{anes}}=0.5$ to conscious sedation. The
advantage of $\delta_{\mathrm{anes}}$ over BIS is that it
extends to children under 5 years, where BIS is unreliable
\citep{davidson2005}.

\subsection{Deafferentation}

In fully deafferented patients (large-fibre sensory neuropathy),
the closed circuit cannot complete its return path
\citep{rothwell1982,cole1995}. The framework predicts that $\Qmo$
collapses to near zero on eye closure even though muscle mass and
motor-neuron integrity are preserved, because the motor command
cannot discharge without its sensory return. The charge subtraction
metric will show $\Qth,\Qpe$ preserved during eyes-open standing and
collapse of $\Qmo$ with eye closure---the inverse pattern from
Parkinson's.

\begin{figure*}[!htbp]
    \centering
    \includegraphics[width=\textwidth]{figures/panel5_clinical.png}
    \caption{\textbf{Clinical Signatures and Run-to-Run Stability.}
    (\textbf{A})~Three-dimensional bar chart showing the mean component charge rates (thought, motor, perception, dream) for four conditions: control, simulated Parkinson's disease, simulated Alzheimer's prodrome, and simulated general anaesthesia. Parkinson's shows a preferential collapse of $\Qmo$ with preserved $\Qth$; Alzheimer's shows preferential $\Qth$ and $\Qpe$ drops with preserved $\Qmo$; anaesthesia drops all cognitive components sharply while motor remains low due to paralysis.
    (\textbf{B})~Bootstrap distribution of $\Qth$ from 200 seed-based resamples of the 86-night record. The point estimate $133.5$~mC/s (red dashed) sits at the centre; the $\pm 1\sigma$ band is $\pm 4.1$~mC/s ($3.1\%$ relative). This narrow distribution satisfies the longitudinal-coverage bound of Theorem~\ref{thm:minlength} at $N=86\gg 30$.
    (\textbf{C})~Scatter of $\Qmo$ versus $\Qth$ for three simulated groups (control, Parkinson's, Alzheimer's). The two disease groups occupy distinct regions: Parkinson's shifts left along the $\Qmo$ axis; Alzheimer's shifts down along $\Qth$. The healthy cloud is centred at $(290, 133)$~mC/s (dashed grey reference lines). The orthogonality of the two disease signatures is the central clinical-use case for the subtraction procedure.
    (\textbf{D})~Empirical cumulative distribution of the anaesthesia depth index $\delta_{\mathrm{anes}}=1-\Qth^{\mathrm{intra}}/\Qth^{\mathrm{pre}}$ for control subjects (green) and simulated general-anaesthesia patients (purple). The two distributions are essentially disjoint: $\delta_{\mathrm{anes}}$ cleanly distinguishes awake from anaesthetised states and, within the anaesthetised group, separates surgical ($\delta\gtrsim 0.8$) from conscious sedation ($\delta\approx 0.5$). The index is computable without EEG, using only consumer-wearable HR and hypnogram streams.}
    \label{fig:clinical}
\end{figure*}


% ═══════════════════════════════════════════════════════════════════
\section{Discussion}
% ═══════════════════════════════════════════════════════════════════

\subsection{What the Dream--Thought Ratio Tells Us}

The observed $\Qdr/\Qth = 0.975$ confirms that thought, defined as
the full cognitive signalling slice, has the same aggregate capacitance
as internally generated REM activity. This is consistent with the
Maquet 1990 and Braun 1997 observation that PET glucose uptake during
REM is $\approx 95\%$ of awake baseline: we here supply the
mechanistic interpretation that the missing $5\%$ is the
REM metabolic multiplier, not a change in underlying capacitive
architecture.

The unity of the ratio is \emph{not} trivially consistent with any
partition of the cognitive slice into separate ``perceptual'' and
``ideational'' components that each have their own capacitance.
Two-capacitance models predict $\Qdr/\Qth<1$ by
more than the $\sqrt{0.95}$ factor whenever the perceptual component
is present during waking but suppressed in REM, because
the active waking circuit has greater aggregate capacitance than
the REM circuit. The observed match to $\sqrt{0.95}$ selects the
single-capacitance, full-cognitive-slice interpretation.

\subsection{Why Consumer Wearables Suffice}

The framework's data requirements are weaker than those of any
quantitative-EEG or metabolic-imaging alternative. A 24-h
heart-rate record, a well-scored hypnogram, and one calibrated
in-clinic ECG suffice to extract all four charge components. For
longitudinal monitoring this is a decisive practical advantage: a
single subject's 86-night record is logistically trivial to collect
but produces a $\pm 3\%$ estimate of $\Qth$. Population-level
epidemiology on $\Qth$ becomes feasible without any study-specific
instrumentation.

\subsection{Limits of Single-Subject Validation}

This paper validates the procedure on one subject. The framework is
single-subject by design---it is meant to deliver individualised
charge phenotypes rather than group statistics---but the clinical
predictions in Section~\ref{sec:clinical} require cohort studies to
establish sensitivity and specificity thresholds. The procedure is
deterministic given the input data streams, so cohort studies reduce
to data aggregation without any additional methodological
development.

\subsection{What the $\kappa$ Scaling Means for Pathology}

The clipping bounds on $f_{\mathrm{perc}}$ at $[0.20,0.60]$ correspond
to the extremes of the cardiac-coupling distribution in healthy
adults. Subjects outside this range (e.g.\ severe autonomic failure
with $\RMSSD\to 0$, or exercise-induced supra-cardiac-coupling with
$\kappa>\kappa_{\mathrm{ref}}$ by more than $3\times$) will show
characteristic $\Qpe$ deviations, which we predict will correlate
with the perceptual-deficit subtypes of dysautonomia
\citep{freeman2011}.

\subsection{Relation to the Closed Non-Grounded Circuit Papers}

Paper~I \citep{sachikonye2026musculo} established the closed
non-grounded circuit framework for the musculoskeletal system and
derived $\Qmo=290$~mC/s from the Hodgkin--Huxley and Kirchhoff axioms.
Paper~II \citep{sachikonye2026rambling} derived the
rambling/trembling decomposition of postural control from the same
circuit. This paper completes the triptych by supplying the
subtraction procedure that separates thought from motor charge in
the intact awake circuit. The unity of the dream--thought ratio is
the framework's strongest quantitative prediction to date.

% ═══════════════════════════════════════════════════════════════════
\section{Conclusions}
% ═══════════════════════════════════════════════════════════════════

\begin{itemize}[leftmargin=1.3em]
  \item Four behavioural conditions span the compartmental charge
  subspace with condition number $6.8$.
  \item A 50/50 split of brain metabolism into housekeeping and
  active signalling, a cardiac-coupling scaling of the perceptual
  sub-component, and the $Q=\sqrt{2CU}$ energy-to-charge map yield
  physiological charge rates.
  \item Single-subject 86-night validation gives
  $\Qth=133.5\pm 4.1$~mC/s, $\Qmo=290.9$~mC/s,
  $\Qpe=70.9$~mC/s, and a dream--thought charge ratio of
  $0.975$, consistent with the framework prediction of
  $\sqrt{0.95}$.
  \item Clinical signatures for Parkinson's, Alzheimer's, anaesthesia,
  and deafferentation follow directly from component-wise charge
  trajectories.
  \item The procedure runs on consumer wearables; cohort studies are
  a matter of data aggregation.
\end{itemize}

\section*{Acknowledgements}

This work is part of a single-subject longitudinal study supported by
the AIMe Registry for Artificial Intelligence and uses data
voluntarily contributed by the author. Sleep records are from an
Oura Ring Gen~3; clinical ECG, PWV, and blood-chemistry data are
from routine annual physicals.

\section*{Data availability}

All data and analysis code used in this paper are available at
\texttt{olduvai-gorge/publications/charge-quantification/}; the
validation runs on \texttt{analyze\_charge.py}.

\bibliographystyle{unsrtnat}
\bibliography{references}

\end{document}
